%! The Nullspace of A: Solving Ax = 0 — Unit 3, Lesson 3.2 — Lesson Plan
\documentclass[10pt]{article}
\usepackage{linalg-article}
\usepackage{linalg-boxes}
\usepackage{graphicx}   % -article does NOT load graphicx; needed for thumbnails/screenshots
\graphicspath{{images/}}
\newcommand{\TallMath}[1]{$\displaystyle #1\rule[-1.4em]{0pt}{3.2em}$}
\newcommand{\vv}{\mathbf{v}}
\newcommand{\ww}{\mathbf{w}}
\newcommand{\xx}{\mathbf{x}}
\newcommand{\yy}{\mathbf{y}}
\newcommand{\bb}{\mathbf{b}}
\newcommand{\svec}{\mathbf{s}}
\newcommand{\zero}{\mathbf{0}}
\newcommand{\UnitNumberName}{Unit 3: The Four Fundamental Subspaces \quad}
\newcommand{\LessonNumberName}{Lesson 3.2: The Nullspace of A: Solving Ax = 0}

\begin{document}
\begin{center}
    {\Large\bfseries \CourseName: \SchoolYear \\
    {\small\bfseries \UnitNumberName \LessonNumberName}}
\end{center}

\begin{tcolorbox}[colback=blush, colframe=burgundy, boxrule=0.9pt, arc=2mm,
    left=2mm, right=2mm, top=1.5mm, bottom=1.5mm]
    \textbf{Primary Objective:} Students can describe the \emph{nullspace} $N(A)$ --- all
    solutions $\xx$ of $A\xx=\zero$ --- and explain why it is a subspace of $\mathbb{R}^n$; solve
    $A\xx=\zero$ by elimination, identifying \emph{pivot columns} and \emph{free columns}; build a
    \emph{special solution} for each free variable (set that variable to $1$, the rest to $0$, solve
    back) and write the whole nullspace as their combinations; and count the free variables as
    $n-r$ to decide whether $N(A)$ is a point, a line, or a plane through the origin.
\end{tcolorbox}

\renewcommand{\arraystretch}{1.4}
\begin{skillbox}[Priority Ideas \& Skills]{goldbox}
\begin{tabularx}{\linewidth}{@{}X|X@{}}
\textbf{Priority Ideas \& Skills} & \textbf{Key Understandings (the ``why'')} \\[2pt]
\hline
\vspace{-6pt}
\begin{itemize}[leftmargin=1.1em, nosep, after=\vspace{2pt}]
  \item Define the \textbf{nullspace} $N(A)=\{\xx:A\xx=\zero\}$ and show it is a subspace of $\mathbb{R}^n$.
  \item Solve $A\xx=\zero$ by \textbf{elimination}; the right side stays $\zero$, so just reduce $A$.
  \item Tell \textbf{pivot columns} from \textbf{free columns}; find one \textbf{special solution} per free variable.
  \item Write $N(A)$ as all combinations of the special solutions; count free variables $=n-r$.
\end{itemize}
&
\vspace{-6pt}
\begin{itemize}[leftmargin=1.1em, nosep, after=\vspace{2pt}]
  \item $N(A)$ collects the \emph{inputs} that produce nothing --- the ``redundancy'' in the columns.
  \item $C(A)$ lives in $\mathbb{R}^m$ (outputs); $N(A)$ lives in $\mathbb{R}^n$ (inputs). Different homes.
  \item A \textbf{free} variable is one you may choose; each choice of $1$ names one direction of $N(A)$.
  \item Free count $n-r$ = the nullspace's dimension: $0\Rightarrow$ point, $1\Rightarrow$ line, $2\Rightarrow$ plane.
\end{itemize}
\end{tabularx}
\end{skillbox}

\begin{skillbox}[{Vocabulary, Concepts, \& Theorems}]{greenbox}
\renewcommand{\arraystretch}{1.3}
\begin{tabularx}{\linewidth}{@{}l X@{}}
\textbf{Nullspace $N(A)$} & All solutions $\xx$ of $A\xx=\zero$; a subspace of $\mathbb{R}^n$ ($n$ = number of columns). \\
\textbf{Pivot column} & A column that holds a pivot after elimination; its variable is a \emph{pivot variable}. \\
\textbf{Free column / free variable} & A column with \emph{no} pivot; its variable is free to choose. \\
\textbf{Special solution} & The nullspace vector from setting one free variable to $1$, the other free variables to $0$, then solving for the pivot variables. \\
\textbf{Rank $r$} & The number of pivots. Free variables $=n-r$. \\
\textbf{$N(A)$ as combinations} & The whole nullspace = all combinations of the special solutions (one per free variable). \\
\end{tabularx}
\end{skillbox}

\begin{fixedskillbox}[Activate Prior Knowledge \& Spiral Review]{blush}
    The Warm-Up rehearses the three tools this lesson runs on:
    \textbf{(1)} checking that a vector solves $A\xx=\zero$ by computing $A\xx$ as a
    \emph{combination of the columns} (Lesson~1.3) --- this pre-verifies the hook's cancelling
    levers; \textbf{(2)} one \emph{elimination} step on a $2\times3$ matrix (Unit~2), noticing that a
    column can end up with \emph{no pivot}; and \textbf{(3)} the Lesson~3.1 \emph{subspace
    requirement}, since $N(A)$ will pass it. Debrief item~1 aloud: a nonzero $\xx$ with $A\xx=\zero$
    is exactly a ``do-nothing'' combination of the columns.
\end{fixedskillbox}

\begin{skillbox}[Hook]{blush}
    A robot arm has two levers; setting them to $\xx=\begin{bsmallmatrix} x_1\\x_2 \end{bsmallmatrix}$
    moves the tip to $A\xx$, a combination of the two lever-direction columns of $A$. In Lesson~3.1
    the question was ``where \emph{can} the tip go?'' ($C(A)$). Today flip it: \textbf{``Which lever
    settings leave the tip exactly where it started --- $A\xx=\zero$?''} If the two levers push along
    the \emph{same} line (columns are multiples), you can push one forward and pull the other back to
    cancel. For $A=\begin{bsmallmatrix} 1 & 2\\ 2 & 4 \end{bsmallmatrix}$, the setting
    $\xx=\begin{bsmallmatrix} -2\\1 \end{bsmallmatrix}$ does nothing:
    $-2\begin{bsmallmatrix} 1\\2 \end{bsmallmatrix}+1\begin{bsmallmatrix} 2\\4 \end{bsmallmatrix}=\zero$.
    The set of \emph{all} such do-nothing settings is the \textbf{nullspace} --- the redundancy hiding
    in the controls.
\end{skillbox}

\begin{skillbox}[Lesson]{blush}
\begin{multicols}{2}
\textbf{1. The nullspace, and why it is a subspace.} The \textbf{nullspace} $N(A)$ is the set of all
solutions $\xx$ of $A\xx=\zero$. It always contains $\zero$ (since $A\zero=\zero$), and it passes the
subspace requirement: if $A\xx_1=\zero$ and $A\xx_2=\zero$ then $A(\xx_1+\xx_2)=\zero+\zero=\zero$ and
$A(c\xx_1)=c\,\zero=\zero$. So $N(A)$ is a subspace of $\mathbb{R}^n$ ($n$ = columns) --- note $C(A)$
lived in $\mathbb{R}^m$, so the two fundamental subspaces have \emph{different} homes.

\textbf{2. Solve $A\xx=\zero$ by elimination.} The right side is all zeros and \emph{stays} zero under
every row operation, so we simply row-reduce $A$. For $A=\begin{bsmallmatrix} 1&1&2\\1&2&3 \end{bsmallmatrix}$,
one step gives $\begin{bsmallmatrix} 1&1&2\\0&1&1 \end{bsmallmatrix}$, then
$\begin{bsmallmatrix} 1&0&1\\0&1&1 \end{bsmallmatrix}$: pivots sit in columns~1 and~2 (\textbf{pivot
columns}, variables $x_1,x_2$); column~3 has \emph{no} pivot (\textbf{free column}, free variable $x_3$).

\columnbreak

\textbf{3. Special solutions.} Give the free variable the value $1$ and solve back: from $x_1+x_3=0$,
$x_2+x_3=0$ with $x_3=1$ we get $x_1=x_2=-1$, so $\svec=\begin{bsmallmatrix} -1\\-1\\1 \end{bsmallmatrix}$.
Every nullspace vector is a multiple of $\svec$. \textbf{Rule:} one special solution \emph{per} free
variable; $N(A)$ = all combinations of them; the number of free variables is $n-r$ ($r$ = pivots).

\textbf{4. Point, line, or plane.} The free count $n-r$ is the nullspace's dimension. \emph{No} free
variables ($r=n$): only $\xx=\zero$, so $N(A)=\{\zero\}$, a \textbf{point}. One free variable: a
\textbf{line} through the origin. Two: a \textbf{plane} through the origin (e.g.
$A=\begin{bsmallmatrix} 1&2&3 \end{bsmallmatrix}$ has $n-r=2$ and $N(A)$ = the plane spanned by
$\begin{bsmallmatrix} -2\\1\\0 \end{bsmallmatrix}$ and $\begin{bsmallmatrix} -3\\0\\1 \end{bsmallmatrix}$).
\end{multicols}
\end{skillbox}

\begin{skillbox}[Active Monitoring]{redbox}
Circulate during the notes practice and Tier work. Watch for and correct:
\begin{itemize}[leftmargin=1.2em, nosep]
  \item saying $N(A)$ lives in $\mathbb{R}^m$ --- it is a set of \emph{inputs} $\xx$, so it lives in $\mathbb{R}^n$ (columns);
  \item choosing pivot columns as ``free'' --- the free columns are exactly the ones \emph{without} a pivot;
  \item forgetting to solve \emph{back} for the pivot variables after setting a free variable to $1$;
  \item writing only the special solution and stopping --- the nullspace is \emph{all multiples/combinations} of it.
\end{itemize}
Cold-call prompts: ``Which columns are free here --- how do you know?'' ``Give me the special solution.''
``Is $N(A)$ a point, a line, or a plane --- and why?'' ``What is $n-r$ for this matrix?''
\end{skillbox}

\begin{skillbox}[Group Work \& Differentiation]{redbox}
\textbf{Find the Nullspace} activity (tiered; mirrors the notes).
\begin{multicols}{3}
\textbf{Tier R --- Remediate.} Verify a given $\xx$ solves $A\xx=\zero$; solve $x_1+2x_2=0$ for the
free-variable form; label pivot vs.\ free columns in an already-reduced matrix.

\columnbreak
\textbf{Tier A --- Approaching.} Reduce a matrix, name the pivot/free columns, build the special
solution(s), describe $N(A)$ as a line or plane, and count $n-r$.

\columnbreak
\textbf{Tier E --- Extension.} Prove $N(A)$ is a subspace in full; connect a nonzero nullspace to
\emph{dependent columns}; and tie $N(A)=\{\zero\}$ to an invertible square matrix / unique solution.
\end{multicols}
\end{skillbox}

\begin{skillbox}[Individual Work \& Assessment]{redbox}
\textbf{Exit Ticket} (independent, no notes): find the nullspace of a small matrix (special solution
+ geometric description); use $n-r$ to state how many free variables a given size/rank forces; and a
\textbf{conceptual/justification check} --- a $3\times3$ matrix with $N(A)=\{\zero\}$: what does that
say about solutions of $A\xx=\bb$? Collect and sort into ``got it / free-vs-pivot slip / stopped at
one special solution'' to decide whether 3.3 opens with a re-teach.
\end{skillbox}

\begin{skillbox}[Reinforcement \& Extension]{goldbox}
\textbf{Homework:} find nullspaces across a battery of matrices (verify membership, reduce, name
free columns, build special solutions, describe the geometry, count $n-r$); contrast where $C(A)$ and
$N(A)$ live. \textbf{Extension:} a nonzero nullspace vector \emph{is} a dependency among the columns;
and $N(A)=\{\zero\}\Leftrightarrow$ a square $A$ is invertible ($A\xx=\bb$ has exactly one solution).
\textbf{Preview:} Lesson~3.3, \emph{The Complete Solution to $A\xx=\bb$} --- one particular solution
plus the whole nullspace as ``wiggle room.''
\end{skillbox}
\end{document}
